%% mtDNA-FM: Supplementary Materials
%% Compile separately or include as appendix

\documentclass[10pt]{article}
\usepackage[margin=2.5cm]{geometry}
\usepackage{amsmath,amssymb,amsthm}
\usepackage{graphicx}
\usepackage{booktabs}
\usepackage{hyperref}
\usepackage{xcolor}

\newcommand{\todo}[1]{\textcolor{red}{\textbf{[TODO: #1]}}}

\begin{document}

\title{\textbf{Supplementary Materials:\\
mtDNA-FM: A Domain-Specialized Foundation Model for Mitochondrial DNA}}
\author{}
\date{}
\maketitle

%% =========================================================================
\section*{S1. Hyperparameters}

\subsection*{S1.1 Pre-training}

\begin{table}[h]
\centering
\caption{Pre-training hyperparameters for Phase~1 and Phase~2.}
\begin{tabular}{lcc}
\toprule
\textbf{Hyperparameter} & \textbf{Phase 1} & \textbf{Phase 2} \\
\midrule
Training steps & 50,000 & 25,000 \\
Warmup steps & 2,000 & 500 \\
Peak learning rate & $1 \times 10^{-4}$ & $5 \times 10^{-5}$ \\
LR schedule & Cosine decay & Cosine decay \\
Batch size (effective) & 256 & 256 \\
Gradient accumulation steps & 8 & 8 \\
Optimizer & AdamW ($\beta_1=0.9, \beta_2=0.999$) & AdamW \\
Weight decay & 0.01 & 0.01 \\
Gradient clipping & 1.0 & 1.0 \\
MLM masking probability & 0.15 & 0.15 \\
Heteroplasmy loss weight ($\lambda_\text{het}$) & 0.0 & 0.3 \\
Checkpoint rotation & Keep last 3 & Keep last 3 \\
Gradient checkpointing & Yes & Yes \\
\bottomrule
\end{tabular}
\end{table}

\subsection*{S1.2 Model Architecture}

\begin{table}[h]
\centering
\caption{mtDNA-FM architecture hyperparameters.}
\begin{tabular}{ll}
\toprule
\textbf{Parameter} & \textbf{Value} \\
\midrule
Hidden size ($d$) & 256 \\
Number of transformer layers & 6 \\
Number of attention heads & 8 \\
Feed-forward intermediate size & 1,024 \\
Attention head dimension & 32 \\
Dropout & 0.1 \\
Attention dropout & 0.1 \\
Max position embeddings & 514 (512 + [CLS] + [SEP]) \\
Vocabulary size & 4,102 \\
Activation function & GELU \\
Layer normalization $\epsilon$ & $10^{-12}$ \\
Total parameters & $\approx$5.8M \\
Positional encoding & Circular sinusoidal (fixed buffer) \\
Genome length for circular PE & 16,569 bp (rCRS) \\
\bottomrule
\end{tabular}
\end{table}

\subsection*{S1.3 Fine-tuning Hyperparameters}

\begin{table}[h]
\centering
\caption{LoRA fine-tuning hyperparameters for each downstream task.}
\begin{tabular}{lccc}
\toprule
\textbf{Hyperparameter} & \textbf{Haplogroup} & \textbf{Pathogenicity} & \textbf{Heteroplasmy} \\
\midrule
LoRA rank ($r$) & 8 & 4 & 4 \\
LoRA $\alpha$ & 16 & 8 & 8 \\
LoRA dropout & 0.05 & 0.05 & 0.05 \\
LoRA target modules & Q, V & Q, V & Q, V \\
Learning rate & $1 \times 10^{-3}$ & $5 \times 10^{-4}$ & $5 \times 10^{-4}$ \\
Epochs & 20 & 30 & 30 \\
Batch size & 32 & 16 & 16 \\
Loss function & Cross-entropy & BCE ($w_+=2.5$) & Huber ($\delta=1.0$) \\
Weight decay & 0.0 & 0.1 & 0.1 \\
Input representation & CLS token & Variant position token & CLS token \\
\bottomrule
\end{tabular}
\end{table}

%% =========================================================================
\section*{S2. Mathematical Derivation: Circular Positional Encoding}

\subsection*{S2.1 Standard sinusoidal PE}

In the original Transformer \citep{vaswani2017attention}, position $p$ is encoded as:
\begin{align}
\text{PE}(p, 2i) &= \sin\!\left(\frac{p}{10000^{2i/d}}\right) \\
\text{PE}(p, 2i+1) &= \cos\!\left(\frac{p}{10000^{2i/d}}\right)
\end{align}
For a sequence of length $L$, $\text{PE}(0, \cdot) \ne \text{PE}(L-1, \cdot)$ in general.
As $d \to \infty$, the angle between $\text{PE}(0)$ and $\text{PE}(L-1)$ grows monotonically
with $L$.

\subsection*{S2.2 Circular modification}

We replace $p$ with the circular angle $\phi_p = 2\pi p / L_\text{genome}$:
\begin{align}
\text{PE}_\text{circ}(p, 2i) &= \sin\!\left(\frac{2\pi p / L_\text{genome}}{10000^{2i/d}}\right) \\
\text{PE}_\text{circ}(p, 2i+1) &= \cos\!\left(\frac{2\pi p / L_\text{genome}}{10000^{2i/d}}\right)
\end{align}

At $p = L_\text{genome}$:
\[
\phi_{L_\text{genome}} = \frac{2\pi L_\text{genome}}{L_\text{genome}} = 2\pi
\]
and $\sin(2\pi) = \sin(0) = 0$, $\cos(2\pi) = \cos(0) = 1$.

Therefore $\text{PE}_\text{circ}(L_\text{genome}, \cdot) = \text{PE}_\text{circ}(0, \cdot)$,
which is the desired circular continuity property.

\subsection*{S2.3 Properties}

\begin{itemize}
\item \textit{Uniqueness:} Positions $p$ and $p'$ with $|p - p'| < L_\text{genome}/2$ receive
distinct encodings (the mapping $p \mapsto \phi_p$ is injective on $[0, L_\text{genome})$).

\item \textit{Smoothness at junction:} $\|\text{PE}_\text{circ}(p) - \text{PE}_\text{circ}(p+1)\|$
is constant across all positions including $p = L_\text{genome} - 1$ (the circular junction),
whereas for standard sinusoidal PE this distance jumps at position 0.

\item \textit{Non-learnable:} Implemented as \texttt{nn.Buffer}, not \texttt{nn.Parameter};
gradients do not flow through the positional encoding.
\end{itemize}

%% =========================================================================
\section*{S3. Per-class Haplogroup Metrics}

\begin{table}[h]
\centering
\caption{Per-class haplogroup classification metrics from fine-tuned \mtdnafm.}
\begin{tabular}{lcccc}
\toprule
\textbf{Haplogroup} & \textbf{Precision} & \textbf{Recall} & \textbf{F1} & \textbf{N test} \\
\midrule
A & \todo{X} & \todo{X} & \todo{X} & \todo{X} \\
B & \todo{X} & \todo{X} & \todo{X} & \todo{X} \\
C & \todo{X} & \todo{X} & \todo{X} & \todo{X} \\
D & \todo{X} & \todo{X} & \todo{X} & \todo{X} \\
E & \todo{X} & \todo{X} & \todo{X} & \todo{X} \\
F & \todo{X} & \todo{X} & \todo{X} & \todo{X} \\
G & \todo{X} & \todo{X} & \todo{X} & \todo{X} \\
H & \todo{X} & \todo{X} & \todo{X} & \todo{X} \\
HV & \todo{X} & \todo{X} & \todo{X} & \todo{X} \\
I & \todo{X} & \todo{X} & \todo{X} & \todo{X} \\
J & \todo{X} & \todo{X} & \todo{X} & \todo{X} \\
K & \todo{X} & \todo{X} & \todo{X} & \todo{X} \\
L0 & \todo{X} & \todo{X} & \todo{X} & \todo{X} \\
L1 & \todo{X} & \todo{X} & \todo{X} & \todo{X} \\
L2 & \todo{X} & \todo{X} & \todo{X} & \todo{X} \\
L3 & \todo{X} & \todo{X} & \todo{X} & \todo{X} \\
L4 & \todo{X} & \todo{X} & \todo{X} & \todo{X} \\
M & \todo{X} & \todo{X} & \todo{X} & \todo{X} \\
N & \todo{X} & \todo{X} & \todo{X} & \todo{X} \\
O & \todo{X} & \todo{X} & \todo{X} & \todo{X} \\
R & \todo{X} & \todo{X} & \todo{X} & \todo{X} \\
T & \todo{X} & \todo{X} & \todo{X} & \todo{X} \\
U & \todo{X} & \todo{X} & \todo{X} & \todo{X} \\
V & \todo{X} & \todo{X} & \todo{X} & \todo{X} \\
W & \todo{X} & \todo{X} & \todo{X} & \todo{X} \\
X & \todo{X} & \todo{X} & \todo{X} & \todo{X} \\
\midrule
\textbf{Macro average} & \todo{X} & \todo{X} & \todo{X} & \todo{X} \\
\bottomrule
\end{tabular}
\label{tab:per_class_haplogroup}
\end{table}

%% =========================================================================
\section*{S4. Dataset Summary}

\begin{table}[h]
\centering
\caption{Datasets used in this study.}
\begin{tabular}{lllll}
\toprule
\textbf{Dataset} & \textbf{Source} & \textbf{Size} & \textbf{Purpose} & \textbf{Version/Date} \\
\midrule
HmtDB & \citet{clima2017hmtdb} & 47,000 seqs & Phase 2 pre-train, fine-tune & Retrieved 2024 \\
NCBI vertebrate mtDNA & NCBI Entrez & 117,615 seqs & Phase 1 pre-train & Retrieved 2024 \\
PhyloTree & \citet{vanoven2009phylotree} & 26 major groups & Haplogroup labels & Build 17 \\
gnomAD v3.1 chrM & \citet{karczewski2020gnomad} & $\sim$5,000 vars & Benign variants & v3.1 \\
ClinVar & \citet{landrum2018clinvar} & $\sim$2,000 vars & Pathogenic variants & 2024-01 \\
MITOMAP & \citet{lott2013mitomap} & $\sim$X vars & External validation & Accessed 2024 \\
HelixMTdb & \citet{bolze2019helixmtdb} & $\sim$X vars & External validation & v1 \\
\bottomrule
\end{tabular}
\end{table}

%% =========================================================================
\section*{S5. DVC Pipeline}

The complete analysis pipeline is managed with DVC. The following stages are defined:

\begin{enumerate}
\item \texttt{download\_hmtdb}: Download and cache HmtDB sequences
\item \texttt{download\_ncbi}: Download vertebrate mtDNA from NCBI Entrez
\item \texttt{download\_variants}: Download gnomAD, ClinVar, PhyloTree variant data
\item \texttt{preprocess}: Quality filter, normalize to rCRS length, train/val/test split
\item \texttt{build\_vocabulary}: Construct 4,102-token k-mer vocabulary
\item \texttt{pretrain\_phase1}: Phase 1 cross-species pre-training (50k steps)
\item \texttt{pretrain\_phase2}: Phase 2 human specialization (25k steps)
\item \texttt{finetune\_haplogroup}: LoRA fine-tuning for haplogroup classification
\item \texttt{finetune\_pathogenicity}: LoRA fine-tuning for pathogenicity prediction
\item \texttt{finetune\_heteroplasmy}: LoRA fine-tuning for heteroplasmy regression
\item \texttt{evaluate}: Compute all metrics, generate evaluation reports
\end{enumerate}

Reproduce the full pipeline: \texttt{dvc repro}

%% =========================================================================
\section*{S6. Ablation Learning Curves}

\todo{Insert ablation training curves here after running ablation scripts.
Figures: (a) Phase 1 MLM loss: circular vs sinusoidal vs learnable PE;
(b) Phase 1 MLM loss: two-phase vs single-phase; (c) Phase 2 MLM loss: with vs without het channel.}

%% =========================================================================
\bibliographystyle{plainnat}
\bibliography{references}

\end{document}
